
\documentclass{article}
\usepackage{colortbl}
\usepackage{makecell}
\usepackage{multirow}
\usepackage{supertabular}

\begin{document}

\newcounter{utterance}

\centering \large Interaction Transcript for game `codenames', experiment `risk\_low', episode 5 with reflexion{-}Qwen3.5{-}9B.
\vspace{24pt}

{ \footnotesize  \setcounter{utterance}{1}
\setlength{\tabcolsep}{0pt}
\begin{supertabular}{c@{$\;$}|p{.15\linewidth}@{}p{.15\linewidth}p{.15\linewidth}p{.15\linewidth}p{.15\linewidth}p{.15\linewidth}}
    \# & $\;$A & \multicolumn{4}{c}{Game Master} & $\;\:$B\\
    \hline

    \theutterance \stepcounter{utterance}  
    & & & \multicolumn{2}{p{0.3\linewidth}}{
        \cellcolor[rgb]{0.9,0.9,0.9}{
            \makecell[{{p{\linewidth}}}]{
                \texttt{\tiny{[GM$|$GM]}}
                {'hidden': {'team': ['phoenix', 'plate'], 'innocent': ['heart', 'plastic', 'war', 'ninja', 'millionaire', 'snowman', 'orange'], 'opponent': [], 'assassin': []}, 'revealed': {'team': {'team': [], 'innocent': [], 'opponent': [], 'assassin': []}, 'opponent': {'team': [], 'innocent': [], 'opponent': [], 'assassin': []}}}
            }
        }
    }
    & & \\ \\

    \theutterance \stepcounter{utterance}  
    & & \multicolumn{4}{p{0.6\linewidth}}{
        \cellcolor[rgb]{0.9,0.9,0.9}{
            \makecell[{{p{\linewidth}}}]{
                \texttt{\tiny{[P1$\langle$GM]}}
                \texttt{You have attempted this same task before and did not succeed. These are the plans you wrote after those attempts.} \\
\texttt{Trial \#1: New plan: I failed because I included "phoenix" in the TARGETS list for the clue "school," which is semantically incorrect and likely violated the rule that the clue must be related to the target words. Additionally, I did not verify that the clue word itself was not a distractor or opponent word before submitting. Next time, I will strictly select a clue word that has a clear, direct semantic link to all listed targets (e.g., "education" for "teacher") and double{-}check that the clue word is not present in the opponent's list, distractor list, or kill words before formatting the response. I will also ensure the TARGETS list only contains words from my remaining team list that are actually related to the clue.} \\
\\ 
\texttt{Trial \#2: New plan: I failed because I provided a clue ("fire") that only related to one remaining word ("phoenix") while ignoring the other ("plate"), allowing the opponent to guess my remaining word on their next turn. The strategy of targeting only one word per turn was too slow. Next time, I will select a single clue word that semantically connects to *all* of my remaining team words (e.g., "material" for "plate" and "phoenix" if a link exists, or finding a broader concept) to clear multiple targets in one move, ensuring I win before the opponent can guess the rest.} \\
\\ 
\texttt{Trial \#3: New plan: I failed because the clue "material" was too broad and ambiguous; "phoenix" is a mythical creature, not a physical material, so my partner likely could not guess it, allowing the opponent to guess my remaining word ("plate") and end the game. Next time, I will select a clue word that has a direct, unambiguous semantic link to *all* listed targets. Since "phoenix" and "plate" share no strong common noun, I should have either chosen a different set of words that share a clear concept (e.g., "fire" for "phoenix" and "burn" for "plate" if both were targets) or, if forced to guess both, I should have realized the impossibility and prioritized guessing the easier word first to secure a win, rather than risking a failed clue that gives the opponent a turn. However, strictly adhering to the rule of one clue for multiple targets, I must ensure the connection is undeniable; if no such word exists for the current pair, I should have only listed the one word I could confidently link (e.g., "metal" for "plate") to avoid losing the turn.} \\
\\ 
\\ 
\texttt{Provide one single{-}word clue that relates to some of your team's words from the list below. You can choose to how many and to which words your clue relates to.} \\
\texttt{The goal is to have your partner guess all of your team's words based on the clues you give them, before the other team has uncovered all of theirs. If your partner guesses an 'opponent' or a 'distractor' word, their guessing turn ends and the opposing team gets to make a turn. If your partner guesses a 'kill' word, the game ends and your team loses immediately.} \\
\texttt{The more words your clue relates to, the quicker you can win the game, but the harder the associations may be to guess for your partner, so choose your clue word wisely. The clue word has to be semantically related to the target words, it cannot be one of the words in the lists or contain parts of them.} \\
\\ 
\texttt{Always give your single{-}word clue and your comma{-}separated list of related target words in the following format and make your answers as short as possible, never include any other text than is required in this form:} \\
\\ 
\texttt{CLUE: <WORD>} \\
\texttt{TARGETS: <TARGETS>} \\
\\ 
\texttt{Your team words are: phoenix, plate.} \\
\texttt{Your opponent's team words are: .} \\
\texttt{Distractor words are: heart, plastic, war, ninja, millionaire, snowman, orange.} \\
\texttt{The kill words are: .} \\
            }
        }
    }
    & \\ \\

    \theutterance \stepcounter{utterance}  
    & \multicolumn{4}{p{0.6\linewidth}}{
        \cellcolor[rgb]{0.9,0.9,0.9}{
            \makecell[{{p{\linewidth}}}]{
                \texttt{\tiny{[P1$\rangle$GM]}}
                \texttt{CLUE: fire} \\
\texttt{TARGETS: phoenix} \\
            }
        }
    }
    & & \\ \\

    \theutterance \stepcounter{utterance}  
    & & & \multicolumn{2}{p{0.3\linewidth}}{
        \cellcolor[rgb]{0.9,0.9,0.9}{
            \makecell[{{p{\linewidth}}}]{
                \texttt{\tiny{[GM$|$GM]}}
                \texttt{fire} \\
            }
        }
    }
    & & \\ \\

    \theutterance \stepcounter{utterance}  
    & & & \multicolumn{2}{p{0.3\linewidth}}{
        \cellcolor[rgb]{0.9,0.9,0.9}{
            \makecell[{{p{\linewidth}}}]{
                \texttt{\tiny{[GM$|$GM]}}
                ['phoenix']
            }
        }
    }
    & & \\ \\

    \theutterance \stepcounter{utterance}  
    & & & \multicolumn{2}{p{0.3\linewidth}}{
        \cellcolor[rgb]{0.9,0.9,0.9}{
            \makecell[{{p{\linewidth}}}]{
                \texttt{\tiny{[GM$|$GM]}}
                {'word': 'phoenix', 'assignment': 'team'}
            }
        }
    }
    & & \\ \\

    \theutterance \stepcounter{utterance}  
    & & & \multicolumn{2}{p{0.3\linewidth}}{
        \cellcolor[rgb]{0.9,0.9,0.9}{
            \makecell[{{p{\linewidth}}}]{
                \texttt{\tiny{[GM$|$GM]}}
                \texttt{opponent has won} \\
            }
        }
    }
    & & \\ \\

    \theutterance \stepcounter{utterance}  
    & \multicolumn{4}{p{0.6\linewidth}}{
        \cellcolor[rgb]{0.95,0.95,0.95}{
            \makecell[{{p{\linewidth}}}]{
                \texttt{\tiny{[P2$\rightarrow$P1]}}
                \texttt{Post{-}episode Reflexion (not sent to GM):} \\
\texttt{New plan: I failed because I chose "fire" as the clue for "phoenix," which is a strong link, but I completely ignored "plate," leaving it vulnerable to the opponent. The strategy of trying to cover all words with a single clue failed because no single word semantically connects both "phoenix" and "plate" without being ambiguous or incorrect. Next time, I will prioritize speed by selecting the most obvious, unambiguous link for the easiest word first (e.g., "metal" for "plate") to secure a guess immediately, even if it means leaving the other word for a subsequent turn, rather than risking a weak or non{-}existent connection for multiple targets.} \\
            }
        }
    }
    & & \\ \\

\end{supertabular}
}

\end{document}
